% SOURCES: docs/0.2.0v/plan.md (§13 Testing Strategy, §14 Performance Budgets);
%          docs/1.0.0v/supervisor-report.md (§7, §10); README.md (Testing);
%          docs/1.0.0v/perf-capture-checklist.md; CHANGELOG.md
% The suite/gate figures are verified for v1.0.0. The render-layer figures are NOT
% yet captured (issue #95) and are reported honestly as pending — do not invent.

\chapter{Testing and Evaluation}
\label{ch:testing}

\section{Testing Strategy}
Testing is automated and layered, mirroring the source tree so that every module is
exercised close to where it is defined. As of v1.0.0 the suite comprises
\textbf{1042 tests across 110 files}, covering the Document Model, the generator,
the runtime snippets, SEO, the export pipeline, the templates, the stores, the user
interface, and the three v1.0.0 surfaces. Table~\ref{tab:suites} summarises the
suites and what each covers.

\begin{table}[H]
  \centering
  \caption{Test suites and their coverage.}
  \label{tab:suites}
  \small
  \begin{tabularx}{\linewidth}{@{}l X@{}}
    \toprule
    \textbf{Suite} & \textbf{Coverage} \\
    \midrule
    \texttt{tests/document/}  & Types, schemas, operations, tokens, validation, presets, migrations \\
    \texttt{tests/generator/} & HTML/CSS/JS emitters, determinism, prettier formatting, invariants \\
    \texttt{tests/runtime/}   & Runtime snippet behaviour in \texttt{jsdom} \\
    \texttt{tests/seo/}       & head / Open Graph / JSON-LD / sitemap / robots and the axe gate \\
    \texttt{tests/export/}    & Full pipeline including minify, dry-run, and self-host fonts \\
    \texttt{tests/templates/} & Blank / portfolio / landing / resume round-trips and axe-cleanliness \\
    \texttt{tests/store/}, \texttt{tests/main/}, \texttt{tests/ui/} & Stores, IPC handlers, renderer components \\
    \texttt{tests/a11y/}      & End-to-end \texttt{axe-core} on rendered output \\
    \texttt{tests/draw/}, \texttt{tests/match/}, \texttt{tests/mcp/} & The v1.0.0 authoring surfaces \\
    \bottomrule
  \end{tabularx}
\end{table}

\section{Quality Gates}
Three gates run on every push and again in continuous integration: linting
(ESLint + Prettier), strict type-checking (the main/preload and the renderer
\texttt{tsconfig}s), and the full test suite. Beyond the three, additional
automated checks enforce properties that ordinary unit tests do not naturally
capture: output determinism (the same document must serialise to identical bytes),
the no-\texttt{position:absolute} invariant, and the \texttt{axe-core} accessibility
gate, which blocks export on any critical or serious violation. Because these gates
are enforced both locally (a pre-push hook) and in CI, a red build cannot reach
\texttt{main}.

\section{Performance Evaluation}
% SOURCES: perf-capture-checklist.md; supervisor-report §7; plan §14
Performance budgets fall into two categories, and the project reports them
separately and honestly. \textbf{Data-layer and export budgets} are measured
headlessly in a dedicated performance lane and pass with wide margins
(Table~\ref{tab:perf}). \textbf{Render-layer budgets} --- drag frame rate, theme
toggle latency, breakpoint switch, layers-tree scroll, and cold-start --- require
the running application and an interactive profiler; at the time of writing these
are \emph{not yet captured}, and are tracked as a post-release task (issue~\#95).
They are deliberately \emph{not} reported here as measured, because doing so would
misrepresent an empirical result that has not been taken; capturing them on the
demo hardware is the top outstanding item (Chapter~\ref{ch:results}).

\begin{table}[H]
  \centering
  \caption{Headless performance results (v1.0.0 performance lane). Re-run and confirm on the target machine before final submission.}
  \label{tab:perf}
  \small
  \begin{tabularx}{\linewidth}{@{}X l l c@{}}
    \toprule
    \textbf{Metric} & \textbf{Budget} & \textbf{Measured} & \textbf{Pass} \\
    \midrule
    Project save (\texttt{.dtw} serialise, 500 elements)  & $<$ 500 ms  & $\approx$ 0.1 ms & \checkmark \\
    Project open (parse + migrate + validate)             & $<$ 1500 ms & $\approx$ 1.3 ms & \checkmark \\
    Undo/redo round-trip (500 elements)                   & $<$ 16 ms   & $\approx$ 0.1 ms & \checkmark \\
    Dry-run / code preview                                & $<$ 500 ms  & $\approx$ 72 ms  & \checkmark \\
    Full export (9 stages incl.\ axe + minify)            & $<$ 10 s    & $\approx$ 0.2--0.3 s & \checkmark \\
    \bottomrule
  \end{tabularx}
\end{table}

\section{Accessibility and Output-Quality Evaluation}
% SOURCES: supervisor-report §5–§6; plan §14; verify:exports (project harness)
Output quality is evaluated by the project's export-verification harness rather
than by inspection. The exported templates are reported to score a Lighthouse SEO
of 1.00 with valid JSON-LD, and the accessibility gate holds critical and serious
violations at zero across every export path (the same gate runs for the interactive
editor and for the headless MCP server). As with the render-layer numbers, the
exact Lighthouse Performance and byte-size figures should be re-captured on the
target machine and recorded before final submission; the harness and the
performance-capture checklist automate that measurement.
